\section{AMC 10 10A 2015}
        \begin{enumerate}
        	\item (\textit{AMC 10 10A 2015}) What is the value of $(2^0-1+5^2+0)^{-1}\times5?$


 $\textbf{(A)}\ -125\qquad\textbf{(B)}\ -120\qquad\textbf{(C)}\ \frac{1}{5}\qquad\textbf{(D)}\ \frac{5}{24}\qquad\textbf{(E)}\ 25$


 
 \textbf{Answer: }C



\textbf{Solution}

$(2^0-1+5^2-0)^{-1}\times5 = (1-1+25-0)^{-1} \times 5 = 25^{-1} \times 5 = \frac{1}{25} \times 5 = \boxed{\textbf{(C) } \, \frac{1}{5}}$.

 


	\item (\textit{AMC 10 10A 2015}) A box contains a collection of triangular and square tiles. There are $25$ tiles in the box, containing $84$ edges total. How many square tiles are there in the box?


 $\textbf{(A)}\ 3\qquad\textbf{(B)}\ 5\qquad\textbf{(C)}\ 7\qquad\textbf{(D)}\ 9\qquad\textbf{(E)}\ 11$


 
 \textbf{Answer: }D



\textbf{Solution}

Let $a$ be the amount of triangular tiles and $b$ be the amount of square tiles.

 Triangles have $3$ edges and squares have $4$ edges, so we have a system of equations.

 We have $a + b$ tiles total, so $a + b = 25$.

 We have $3a + 4b$ edges total, so $3a + 4b = 84$.

 Multiplying the first equation by $3$ on both sides gives $3a + 3b = 3(25) = 75$.

 Second equation minus the first equation gives $b = 9$, so the answer is $\boxed{\textbf{(D) }9}$.

 



\textbf{Solution 2}

If all of the tiles were triangles, there would be $75$ edges. This is not enough, so there needs to be some squares. Trading a triangle for a square results in one additional edge each time, so we must trade out $9$ triangles for squares. Answer: $\boxed{\textbf{(D) }9}$

 



\textbf{Solution 3}

Let $x$ be the number of square tiles. A square has $4$ edges, so the total number of edges from the square tiles is $4x$. There are $25$ total tiles, which means that there are $25-x$ triangle tiles. A triangle has $3$ edges, so the total number of edges from the triangle tiles is $3(25-x)$. Together, the total number of edges is $4x+3(25-x)=84$. Solving our equation, we get that $x=9$ which means that our answer is $\boxed{\textbf{(D) }9}$.

 


	\item (\textit{AMC 10 10A 2015}) Ann made a 3-step staircase using 18 toothpicks as shown in the figure.  How many toothpicks does she need to add to complete a 5-step staircase?


 
                    \begin{figure}[H]
                        \centering
                        \includegraphics[width=0.5\textwidth]{Problem/AMC_10_10A_2015_3_Problem_Figure_1.png}
                    \end{figure}
                    


 $\textbf{(A)}\ 9\qquad\textbf{(B)}\ 18\qquad\textbf{(C)}\ 20\qquad\textbf{(D)}\ 22\qquad\textbf{(E)}\ 24$


 
 \textbf{Answer: }D



\textbf{Solution}

We can see that a $1$-step staircase requires $4$ toothpicks and a $2$-step staircase requires $10$ toothpicks. Thus, to go from a $1$-step to $2$-step staircase, $6$ additional toothpicks are needed and to go from a $2$-step to $3$-step staircase, $8$ additional toothpicks are needed. Applying this pattern, to go from a $3$-step to $4$-step staircase, $10$ additional toothpicks are needed and to go from a $4$-step to $5$-step staircase, $12$ additional toothpicks are needed. Our answer is $10+12=\boxed{\textbf{(D)}\ 22}$

 



\textbf{Solution 2}

Alternatively, we can see with the $3$-step staircase has $2[2(3)+2+1]=18$ toothpicks. Generalizing, we see that a staircase with $x$ steps has $2[2x+(x-1)+(x-2)+...+1]$ toothpicks. So, for $x=5$ steps, we have $2[2(5)+4+3+2+1]=40$ toothpicks. So our answer is $40-18=22$ or $D$.

 



\textbf{Solution 3}

If one is too lazy to derive a formula for the number of picks needed for a given number of steps, one can simply see that to get to $4$ steps, we add two blobs that have three picks each (the top and the right), and two more blobs that have two blocks each to form the steps. This adds $2\cdot3+2\cdot2=10$ picks. Then, to get to $5$ steps, we add two more edge blobs with $3$ picks each and $3$ more blobs that have two picks each. We add $2\cdot3+3\cdot2=12$ more for a total increase of $10+12=\boxed{\textbf{(D)}~22}.$ 

 


	\item (\textit{AMC 10 10A 2015}) Pablo, Sofia, and Mia got some candy eggs at a party.  Pablo had three times as many eggs as Sofia, and Sofia had twice as many eggs as Mia.  Pablo decides to give some of his eggs to Sofia and Mia so that all three will have the same number of eggs.  What fraction of his eggs should Pablo give to Sofia?


 $\textbf{(A)}\ \frac{1}{12}\qquad\textbf{(B)}\ \frac{1}{6}\qquad\textbf{(C)}\ \frac{1}{4}\qquad\textbf{(D)}\ \frac{1}{3}\qquad\textbf{(E)}\ \frac{1}{2}$


 
 \textbf{Answer: }B



\textbf{Solution 1}

Assign a variable to the number of eggs Mia has, say $m$.  Then, because we are given that Sofia has twice the number of eggs Mia has, Sofia has $2m$ eggs, and Pablo, having three times the number of eggs as Sofia, has $6m$ eggs.

 For them to all have the same number of eggs, they must each have $\frac{m+2m+6m}{3} = 3m$ eggs.  This means Pablo must give $2m$ eggs to Mia and $m$ eggs to Sofia, so the answer is $\frac{m}{6m} = \boxed{\textbf{(B) }\frac{1}{6}}$

 



\textbf{Solution 2}

Let Mia have $1$ egg. Sofia and Pablo would then have $2$ and $6$ eggs, respectively. If we split the eggs evenly, each would have $\frac{1 + 2 + 6}{3} = \frac{9}{3} = 3$ eggs. Since Sofia needs $3 - 2 = 1$ more egg, Pablo must give her $1$ of his $6$ eggs, or $\boxed{\textbf{(B) }\frac{1}{6}}$

 


	\item (\textit{AMC 10 10A 2015}) Mr. Patrick teaches math to $15$ students. He was grading tests and found that when he graded everyone's test except Payton's, the average grade for the class was $80$. After he graded Payton's test, the test average became $81$. What was Payton's score on the test?


 $\textbf{(A)}\ 81\qquad\textbf{(B)}\ 85\qquad\textbf{(C)}\ 91\qquad\textbf{(D)}\ 94\qquad\textbf{(E)}\ 95$


 
 \textbf{Answer: }E



\textbf{Solution}

If the average of the first $14$ peoples' scores was $80$, then the sum of all of their tests is $14 \cdot 80 = 1120$.  When Payton's score was added, the sum of all of the scores became $15 \cdot 81 = 1215$.  So, Payton's score must be $1215-1120 = \boxed{\textbf{(E) }95}$

 \textbf{Alternate Solution} The average of a set of numbers is the value we get if we evenly distribute the total across all entries. So assume that the first $14$ students each scored $80$. If Payton also scored an $80$, the average would still be $80$. In order to increase the overall average to $81$, we need to add one more point to all of the scores, including Payton's. This means we need to add a total of $15$ more points, so Payton needs $80+15 = \boxed{\textbf{(E) }95}$

 


	\item (\textit{AMC 10 10A 2015}) The sum of two positive numbers is $5$ times their difference. What is the ratio of the larger number to the smaller number?


 $\textbf{(A)}\ \frac{5}{4}\qquad\textbf{(B)}\ \frac{3}{2}\qquad\textbf{(C)}\ \frac{9}{5}\qquad\textbf{(D)}\ 2 \qquad\textbf{(E)}\ \frac{5}{2}$


 
 \textbf{Answer: }B



\textbf{Solution 1}

Let $a$ be the bigger number and $b$ be the smaller.

 $a + b = 5(a - b)$.

 Multiplying out gives $a + b = 5a - 5b$ and rearranging gives $4a = 6b$ and factorised into $2a = 3b$ and then solving gives

 $\frac{a}{b} = \frac32$, so the answer is $\boxed{\textbf{(B) }\frac32}$.

 



\textbf{Solution 2}

Without loss of generality, let the two numbers be $3$ and $2$, as they clearly satisfy the condition of the problem. The ratio of the larger to the smaller is $\boxed{\textbf{(B) }\frac32}$.

 


	\item (\textit{AMC 10 10A 2015}) How many terms are there in the arithmetic sequence $13$, $16$, $19$, . . ., $70$, $73$?


 $\textbf{(A)}\ 20\qquad\textbf{(B)}\ 21\qquad\textbf{(C)}\ 24\qquad\textbf{(D)}\ 60\qquad\textbf{(E)}\ 61$


 
 \textbf{Answer: }B



\textbf{Solution 1}

$73-13 = 60$, so the amount of terms in the sequence $13$, $16$, $19$, $\dotsc$, $70$, $73$ is the same as in the sequence $0$, $3$, $6$, $\dotsc$, $57$, $60$. 

 In this sequence, the terms are the multiples of $3$ going up to $60$, and there are $20$ multiples of $3$ in $60$. 

 However, the number 0 must also be included, adding another multiple. So, the answer is $\boxed{\textbf{(B)}\ 21}$.

 



\textbf{Solution 2}

Using the formula for arithmetic sequence's nth term, we see that $a + (n-1)d \Longrightarrow13 + (n-1)3 =73, \Longrightarrow n = 21$$\boxed{\textbf{(B)}\ 21}$.

 



\textbf{Solution 3}

Minus each of the terms by $12$ to make the the sequence $1 , 4 , 7,..., 61$. $\frac{61-1}{3}=20, 20 + 1 = 21$

 $\boxed{\textbf{(B)}\ 21}$.

 



\textbf{Solution 4}

Subtract each of the terms by $10$ to make the sequence $3 , 6 , 9,..., 60, 63$. Then divide the each term in the sequence by $3$ to get $1, 2, 3,..., 20, 21$. Now it is clear to see that there are $21$ terms in the sequence.$\boxed{\textbf{(B)}\ 21}$.

 


	\item (\textit{AMC 10 10A 2015}) Two years ago Pete was three times as old as his cousin Claire. 2 years before that, Pete was four times as old as Claire. In how many years will the ratio of their ages be $2$ : $1$?


 $\textbf{(A)}\ 2\qquad\textbf{(B)}\ 4\qquad\textbf{(C)}\ 5\qquad\textbf{(D)}\ 6\qquad\textbf{(E)}\ 8$


 
 \textbf{Answer: }B



\textbf{Solution}

This problem can be converted to a system of equations. Let $p$ be Pete's current age and $c$ be Claire's current age. 

 The first statement can be written as $p-2=3(c-2)$. The second statement can be written as $p-4=4(c-4)$

 To solve the system of equations:

 
$p=3c-4$

 $p=4c-12$

 $3c-4=4c-12$

 $c=8$

 $p=20.$

 Let $x$ be the number of years until Pete is twice as old as his cousin.

 $20+x=2(8+x)$

 $20+x=16+2x$

 $x=4$

 The answer is $\boxed{\textbf{(B) }4}$.

 


	\item (\textit{AMC 10 10A 2015}) Two right circular cylinders have the same volume. The radius of the second cylinder is $10\%$ more than the radius of the first. What is the relationship between the heights of the two cylinders?


 $\textbf{(A)}\  \text{The second height is } 10\% \text{ less than the first.}$
$\textbf{(B)}\  \text{The first height is } 10\% \text{ more than the second.}$
$\textbf{(C)}\ \text{The second height is } 21\% \text{ less than the first.}$
$\textbf{(D)}\ \text{The first height is } 21\% \text{ more than the second.}$
$\textbf{(E)}\ \text{The second height is } 80\% \text{ of the first.}$


 
 \textbf{Answer: }D



\textbf{Solution}

Let the radius of the first cylinder be $r_1$ and the radius of the second cylinder be $r_2$. Also, let the height of the first cylinder be $h_1$ and the height of the second cylinder be $h_2$. We are told 
 \[r_2=\frac{11r_1}{10}\] 
 \[\pi r_1^2h_1=\pi r_2^2h_2\] Substituting the first equation into the second and dividing both sides by $\pi$, we get 
 \[r_1^2h_1=\frac{121r_1^2}{100}h_2\implies h_1=\frac{121h_2}{100}.\] Therefore, $\boxed{\textbf{(D)}\ \text{The first height is } 21\% \text{ more than the second.}}$

 


	\item (\textit{AMC 10 10A 2015}) How many rearrangements of $abcd$ are there in which no two adjacent letters are also adjacent letters in the alphabet?  For example, no such rearrangements could include either $ab$ or $ba$.


 $\textbf{(A)}\ 0\qquad\textbf{(B)}\ 1\qquad\textbf{(C)}\ 2\qquad\textbf{(D)}\ 3\qquad\textbf{(E)}\ 4$


 
 \textbf{Answer: }C



\textbf{Solution 1}

The first thing one would want to do is place a possible letter that works and then stem off of it.  For example, if we start with an $a$, we can only place a $c$ or $d$ next to it.  Unfortunately, after that step, we can't do too much, since:

 $acbd$ is not allowed because of the $cb$, and $acdb$ is not allowed because of the $cd$. 

 We get the same problem if we start with a $d$, since a $b$ will have to end up in the middle, causing it to be adjacent to an $a$ or $c$.

 If we start with a $b$, the next letter would have to be a $d$, and since we can put an $a$ next to it and then a $c$ after that, this configuration works.  The same approach applies if we start with a $c$.

 So the solution must be the two solutions that were allowed, one starting from a $b$ and the other with a $c$, giving us:

 
 \[1 + 1 = \boxed{\textbf{(C)}\ 2}.\]

 


 



\textbf{Solution 2 (Casework)}

\textbf{Case 1: the first letter is A}

 \textit{Subcase 1: the second letter is C}

 The next letter must either be B or D, both of which do not satisfy the conditions.

 \textit{Subcase 2: the second letter is D}

 The third letter is forced to be B and the fourth is forced to be C, but this doesn't work because B and C are next to each other.

 \textbf{Case 2: the first letter is B}

 The next letter is forced to be D, the third letter forced to be A and the last forced to be C. This works.

 \textbf{Case 3: the first letter is C}

 The next letter is forced to be A, the third letter is D and the last letter is B. This works.

 \textbf{Case 4: the first letter is D}

 \textit{Subcase 1: the second letter is A}

 The third letter has to be C and the fourth has to be B but this doesn't work because B and C are next to each other.

 \textit{Subcase 2: the second letter is B}

 The third letter cannot be A or C, so this doesn't work.

 Summing the cases, there are two that work: $BDAC$ and $CADB$ $\Longrightarrow \boxed{\textbf{(C)}\ 2}$.

 


 



\textbf{Solution 3}

We note that $a$ must be adjacent to $c$ or $d$, $b$ must be adjacent to $d$, $c$ must be adjacent to $a$, and $d$ must be adjacent to $a$ or $b$. In other words, the way the letters can be connected is 
 \[c-a-d-b.\] Therefore, the only two possible arrangements are $cadb$ and $bdac$, which gives $\boxed{\textbf{(C)}\ 2}$.

 


 



\textbf{Solution 4}

Notice that $b$ and $c$ cannot be in the middle because they would be forced to be next to at least one letter that is next to them in the alphabet(For example, if $b$ was in the middle the only letter that can be next to $b$ is $d$ but since it's in the middle it needs another letter, which means that it's not possible). Therefore you either have a case starting with $b$(then it has to be followed by $d$ then $a$ then $c$) or starting with $c$, amounting to $\boxed{\textbf{(C)}\ 2}$ cases

 


	\item (\textit{AMC 10 10A 2015}) The ratio of the length to the width of a rectangle is $4$ : $3$. If the rectangle has diagonal of length $d$, then the area may be expressed as $kd^2$ for some constant $k$. What is $k$?


 $\textbf{(A)}\ \frac{2}{7}\qquad\textbf{(B)}\ \frac{3}{7}\qquad\textbf{(C)}\ \frac{12}{25}\qquad\textbf{(D)}\ \frac{16}{25}\qquad\textbf{(E)}\ \frac{3}{4}$


 
 \textbf{Answer: }C



\textbf{Solution 1}

Let the rectangle have length $4x$ and width $3x$. Then by $3-4-5$ triangles (or the Pythagorean Theorem), we have $d = 5x$, and so $x = \dfrac{d}{5}$. Hence, the area of the rectangle is $3x \cdot 4x = 12x^2 = \dfrac{12d^2}{25}$, so the answer is $\boxed{\textbf{(C)}\;\frac{12}{25}}$

 



\textbf{Solution 2}

Let the side lengths of the rectangle be $3$ and $4$. Therefore, the diagonal of the rectangle is $5$. The area of the rectangle would be $3\cdot{4}=12.$ Thus, $25k=12$, so $k=\boxed{\text{(C)} \frac{12}{25}}.$

 


	\item (\textit{AMC 10 10A 2015}) Points $(\sqrt{\pi}, a)$ and $(\sqrt{\pi}, b)$ are distinct points on the graph of $y^2+x^4=2x^2y+1$. What is $|a-b|$?


 $\textbf{(A) }1\qquad\textbf{(B) }\dfrac{\pi}{2}\qquad\textbf{(C) }2\qquad\textbf{(D) }\sqrt{1+\pi}\qquad\textbf{(E) }1+\sqrt{\pi}$


 
 \textbf{Answer: }C



\textbf{Solution 1}

Since points on the graph make the equation true, substitute $\sqrt{\pi}$ in to the equation and then solve to find $a$ and $b$.

 $y^2 + \sqrt{\pi}^4 = 2\sqrt{\pi}^2 y + 1$

 $y^2 + \pi^2 = 2\pi y + 1$

 $y^2 - 2\pi y + \pi^2 = 1$

 $(y-\pi)^2 = 1$

 $y-\pi = \pm 1$

 $y = \pi + 1$ 

 $y = \pi - 1$

 There are only two solutions to the equation $(y-\pi)^2 = 1$, so one of them is the value of $a$ and the other is $b$. The order does not matter because of the absolute value sign.

 $| (\pi + 1) - (\pi - 1) | = 2$

 The answer is $\boxed{\textbf{(C) }2}$

 



\textbf{Solution 2}

This solution is very closely related to Solution \#1 and just simplifies the problem earlier to make it easier.

 $y^2 + x^4 = 2x^2 y + 1$ can be written as $x^4-2x^2y+y^2=1$. Recognizing that this is a binomial square, simplify this to $(x^2-y)^2=1$. This gives us two equations:

 $x^2-y=1$ and $x^2-y=-1$.

 One of these $y$'s is $a$ and one is $b$. Substituting $\sqrt{\pi}$ for $x$, we get $a=\pi+1$ and $b=\pi-1$. 

 So, $|a-b|=|(\pi+1)-(\pi-1)|=2$.

 The answer is $\boxed{\textbf{(C) }2}$

 



\textbf{Solution 3}

This solution is similar to Solution \#1 but uses a different way to find $y$ at the end.

 Just like Solution \#1, we arrive at the conclusion that $y^2 - 2\pi y + \pi^2 = 1$.

 Simplifying we get:

 $y^2 - 2\pi y + \pi^2 -1 = 0$

 We now can factor this quadratic. We must find two terms that multiply to $\pi^2 -1$ and add to $2\pi$. 

 These terms are $\pi+1$ and $\pi-1$. 

 Subtracting one from the other, we get $2$.

 Thus, the answer is $\boxed{\textbf{(C) }2}$

 -DuckDuckGooseGoose

 


	\item (\textit{AMC 10 10A 2015}) Claudia has 12 coins, each of which is a 5-cent coin or a 10-cent coin. There are exactly 17 different values that can be obtained as combinations of one or more of her coins. How many 10-cent coins does Claudia have?


 $\textbf{(A) }3\qquad\textbf{(B) }4\qquad\textbf{(C) }5\qquad\textbf{(D) }6\qquad\textbf{(E) }7$


 
 \textbf{Answer: }C



\textbf{Solution 1}

Let Claudia have $x$ 5-cent coins and $\left( 12 - x \right)$ 10-cent coins. It is easily observed that any multiple of $5$ between $5$ and $5x + 10(12 - x) = 120 - 5x$ inclusive can be obtained by a combination of coins. Thus, $24 - x = 17$ combinations can be made, so $x = 7$. But the answer is not $7,$ because we are asked for the number of 10-cent coins, which is $12 - 7 = \boxed{\textbf{(C) } 5}$ ~kurt

 



\textbf{Solution 2}

Since the coins are 5-cent and 10-cent, all possible values that can be made will be multiples of $5.$ To have exactly 17 different multiples of $5,$ we will need to make up to $85$ cents. There are two solutions from here:

 1. Hare and Chicken approach

 We assume that the 12 coins are all 10-cent coins. The sum will now be (12)(10) which is 120. 120 is 35 more than 85, and every 10-cent coin we switch for a 5-cent one will reduce the sum by 10-5=5 cents. Since 35/5=7, we will need to switch 7 10-cent coins for 5-cent ones, meaning we will have 7 5-cent coins and 12-7=5 10-cent coins. Therefore, the correct answer is answer choice $\boxed{\textbf{(C) } 5}$

 
2. Algebra approach

 We assume we have (x) 5-cent coins and (12-x)10-cent coins. Then we have the algebra formula 5x+10(12-x)=85 Simplyfying that, we get 120-5x=85, 5x=35, x=7. And since x is the number of 5-cent coins, the number of 10-cent coins is 12-x or 12-7=5. Therefore, the correct answer is answer choice $\boxed{\textbf{(C) } 5}$

 -Alina

 



\textbf{Solution 3 (Quick Insight)}

Notice that for every $d$ dimes, any multiple of $5$ less than or equal to $10d + 5(12-d)$ is a valid arrangement. Since there are $17$ in our case, we have $10d + 5(12-d)=17 \cdot 5 \Rightarrow d=5$. Therefore, the answer is $\boxed{\textbf{(C) } 5}$.

 



\textbf{Solution 4}

Dividing by 5cents to reduce clutter:

 $D$ double coins and $S$ single coins can reach any value between $1$ and $2D + S$. Set $2D + S = 17$ and $D + S = 12$. 

 Subtract to get $D = \boxed{\textbf{(C) } 5}$.

 


	\item (\textit{AMC 10 10A 2015}) The diagram below shows the circular face of a clock with radius $20$ cm and a circular disk with radius $10$ cm externally tangent to the clock face at $12$ o'clock. The disk has an arrow painted on it, initially pointing in the upward vertical direction. Let the disk roll clockwise around the clock face. At what point on the clock face will the disk be tangent when the arrow is next pointing in the upward vertical direction?


 
                    \begin{figure}[H]
                        \centering
                        \includegraphics[width=0.5\textwidth]{Problem/AMC_10_10A_2015_14_Problem_Figure_1.png}
                    \end{figure}
                    


 $\textbf{(A) }\mathrm{2 o'clock} \qquad\textbf{(B) }\mathrm{3 o'clock} \qquad\textbf{(C) }\mathrm{4 o'clock} \qquad\textbf{(D) }\mathrm{6 o'clock} \qquad\textbf{(E) }\mathrm{8 o'clock}$


 
 \textbf{Answer: }C



\textbf{Solution 1 (Important Clarification)}

The below solutions don't do a great job on telling you why the answer is not $\textbf{(D) }\text{6 o' clock}$. For those who don't understand it, the below video is great on explaining it.

 


	\item (\textit{AMC 10 10A 2015}) Consider the set of all fractions $\frac{x}{y},$ where $x$ and $y$ are relatively prime positive integers. How many of these fractions have the property that if both numerator and denominator are increased by $1$, the value of the fraction is increased by $10\%$?


 $\textbf{(A) }0\qquad\textbf{(B) }1\qquad\textbf{(C) }2\qquad\textbf{(D) }3\qquad\textbf{(E) }\text{infinitely many}$


 
 \textbf{Answer: }B



\textbf{Solution 1}

You can create the equation $\frac{x+1}{y+1}=\frac{11x}{10y}$

 Cross multiplying and combining like terms gives $xy + 11x - 10y = 0$.

 This can be factored into $(x - 10)(y + 11) = -110$.

 $x$ and $y$ must be positive, so $x > 0$ and $y > 0$, so $x - 10> -10$ and $y + 11 > 11$.

 Using the factors of 110, we can get the factor pairs: $(-1, 110),$ $(-2, 55),$ and $(-5, 22).$

 But we can't stop here because $x$ and $y$ must be relatively prime.

 $(-1, 110)$ gives $x = 9$ and $y = 99$. $9$ and $99$ are not relatively prime, so this doesn't work.

 $(-2, 55)$ gives $x = 8$ and $y = 44$. This doesn't work.

 $(-5, 22)$ gives $x = 5$ and $y = 11$. This does work.

 We found one valid solution so the answer is $\boxed{\textbf{(B) }1}$.

 



\textbf{Solution 2}

The condition required is $\frac{x+1}{y+1}=\frac{11}{10}\cdot\frac{x}{y}$.

 Observe that $x+1 > \frac{11}{10}\cdot x$ so $x$ is at most $9.$

 By multiplying by $\frac{y+1}{x}$ and simplifying we can rewrite the condition as $y=\frac{11x}{10-x}$.  Since $x$ and $y$ are integer, this only has solutions for $x\in\{5,8,9\}$.  However, only the first yields a $y$ that is relative prime to $x$.

 There is only one valid solution so the answer is $\boxed{\textbf{(B) }1}$

 



\textbf{Solution 3}

So from this question, we can get \(\frac{x+1}{y+1} = \frac{11x}{10y}\). We can transform this equation into \(x + 11 \cdot \left( \frac{x}{y} \right) = 10\). Two numbers are added to get 10, and one of them, \(x\), is a positive and prime integer. So the other number also has to be a positive integer. Therefore, \(11 \cdot \left( \frac{x}{y} \right)\) is a positive integer. The only possibility of this being true is when \(y\) and 11 cancel out, leaving a singular \(x\). So \(y = 11\) and \(x + x = 10\). Therefore, \(y = 11\) and \(x = 5\).

 


	\item (\textit{AMC 10 10A 2015}) If $y+4 = (x-2)^2, x+4 = (y-2)^2$, and $x \neq y$, what is the value of $x^2+y^2$?


 $\textbf{(A) }10\qquad\textbf{(B) }15\qquad\textbf{(C) }20\qquad\textbf{(D) }25\qquad\textbf{(E) }\text{30}$


 
 \textbf{Answer: }B



\textbf{Solution 1}

Note that we can add the two equations to yield the equation

 $x^2 + y^2 - 4x - 4y + 8 = x + y + 8.$

 Moving terms gives the equation

 $x^2+y^2=5 \left( x + y \right).$

 We can also subtract the two equations to yield the equation

 $x^2 - y^2 - 4x +4y = y - x.$

 Moving terms gives the equation

 $x^2 - y^2 = (x + y)(x - y) = 3x - 3y = 3(x - y).$

 Because $x \neq y,$ we can divide both sides of the equation by $x - y$ to yield the equation

 $x + y = 3.$

 Substituting this into the equation for $x^2 + y^2$ that we derived earlier gives

 $x^2 + y^2 = 5 \left( x + y \right) = 5 \left( 3 \right) = \boxed{\textbf{(B) } 15}$

 Note: This may not be the most efficient solution, but it follows a clear and continuous thought process, and is therefore easy to follow.

 



\textbf{Solution 2 (Algebraic)}

Subtract $4$ from the left hand side of both equations, and use difference of squares to yield the equations

 $x = y(y-4)$  and  $y = x(x-4)$.

 It may save some time to find two solutions, $(0, 0)$ and $(5, 5)$, at this point. However, $x = y$ in these solutions.

 
Substitute $y = x(x-4)$ into $x = y(y-4)$.

 
This gives the equation

 $x = x(x-4)(x^2-4x-4)$

 which can be simplified to

 $x(x^3 - 8x^2 +12x + 15) = 0$.

 Knowing $x = 0$ and $x = 5$ are solutions is now helpful, as you divide both sides by $x(x-5)$. This can also be done using polynomial division to find $x = 5$ as a factor. This gives

 $x^2 - 3x -3 = 0$.

 Because the two equations $x = y(y-4)$ and $y = x(x-4)$ are symmetric, the $x$ and $y$ values are the roots of the equation, which are $x = \frac{3 + \sqrt{21}}{2}$ and $x = \frac{3 - \sqrt{21}}{2}$.

 Squaring these and adding them together gives

 $\frac{3^2 + 21 + 6\sqrt{21}}{4} + \frac{3^2 + 21 - 6\sqrt{21}}{4} = \frac{2(3^2 +21)}{4} = \boxed{\textbf{(B) } 15}$.

 \textbf{Alternative Ending} Note that $x^2+y^2 = (x+y)^2 - 2xy$. Using Vieta's, we find that $(x+y)^2 = 3^2 = 9$ and $-2xy = -2\cdot 3 = -6$. Subtracting them gives $\boxed{\textbf{(B) } 15}$.

 



\textbf{Solution 3}

By graphing the two equations on a piece of graph paper, we can see that the point where they intersect that is not on the line $y=x$ is close to the point $(4,-1)$ (or $(-1, 4)$). $(-1)^2+4^2=17$, and the closest answer choice to $17$ is $\boxed{\textbf{(B) } 15}$.

 \textbf{Note} This is risky, as $20$ could be a viable answer too. Do not use this method unless you're sure about the answer. In other words, this solution is less reliable than the others, so only use it if you can't do the other methods.

 



\textbf{Solution 4 (When you can't algebraically manipulate)}

This, by the way, is a really cheap way of solving the problem but as long as you get an answer, it doesn't matter.

 Looking at the first given equation, begin searching for solutions. Notice how $(4,0)$ works but when plugged into the

 second equation, you get $8=4$. Now, if you decrease the $y$ value by $1$ to obtain $(2+\sqrt3,-1)$, plugging it into 

 the second equation will yield $6+\sqrt3=9$. Now, notice how the LHS is now less than the RHS as opposed to what it had 

 been when we plugged $(4,0)$ into the second given equation. We then conclude that there must be a solution between 

 $y=0$ and $y=-1$, so calculating $x^2+y^2$ of $(4,0)$ and $(2+\sqrt3,-1)$ we obtain

 
 \[4^2+0^2=16\] and
 \[(2+\sqrt3)^2+(-1)^2=8+4\sqrt3\approx14.9\]

 Thus we know the answer to be $\boxed{\textbf{(B) } 15}$.

 This is a really bad solution, but it works. :/

 



\textbf{Solution 5 (Very in depth and reasoned)}

Note: In this solution, I explain my entire thought process and the reason why I did each step, so you can read a more intuitive solution rather than a contrived one.

 We first notice that each equation has a $4$, which inspires us to isolate the $4$ on both equations and equate them to see if we can gather any useful information on $x$ and $y.$ 

 Doing this we get $(x-2)^2-y = (y-2)^2- x.$ Now you might be tempted to expand, but notice how we have two perfect squares in our equation: $(x-2)^2$ and $(y-2)^2,$ which motivates us to take advantage of the difference of squares. For simplicity, we suppose that $(x-2) = a$ and $(y-2) = b.$ Thus our equation is now $a^2-y = b^2-x.$ We manipulate this equation by getting the LHS in a form that lets us do difference of squares, so we can rewrite as $a^2-b^2 = -(x-y).$

 Now we can finally use our difference of squares to get $(a-b)(a+b) =  -(x-y).$

 Let's now substitute $a$ and $b$ back to get information solely for $x$ and $y.$ 

 Doing this we get, $(x-y)(x+y-4) = -(x-y).$

 Now let us divide $x-y$ on both sides and clean up the equation. When we do this, we get that $x+y = 3.$

 We got $x+y = 3,$ which means that this is fact is only of use if we can obtain it in another way, so we add the LHS and the RHS of both equations and equate them to make our fact of essence. 

 Doing that, we get $x+y +8 = (x-2)^2+ (y-2)^2.$ Also, when we set this equation up we not only do we make our fact useful, but we also can see that after expanding the RHS, we get $x^2+y^2$ on the RHS too, which makes it even a better move that we added both equations because remember our goal is to find $x^2+y^2.$

 Expanding and plugging in $x+y = 3$ we get $11 = x^2+y^2-4(x+y)+8.$ $x+y=3,$ so we can plug that in and we get $x^2+y^2 = 15$ or $\boxed{\textbf{(B) } 15}$ after cleaning up some algebra.

 



\textbf{Solution 6 (Quickest)}

Subtract the two equations to get $(x-y)= (x+y-4)(y-x) \iff x+y=3$. Plugging back into any of the original equations yields $x^2-3x=3$. However, we know $x^2+y^2= x^2+(3-x)^2= 2x^2-6x+9$. Then, we can just multiply the first by two and add $9$ to get $\boxed{\textbf{(B) } 15}$

 


	\item (\textit{AMC 10 10A 2015}) A line that passes through the origin intersects both the line $x=1$ and the line $y=1+\frac{\sqrt{3}}{3}x$. The three lines create an equilateral triangle. What is the perimeter of the triangle?


 $\textbf{(A) }2\sqrt{6}\qquad\textbf{(B) }2+2\sqrt{3}\qquad\textbf{(C) }6\qquad\textbf{(D) }3+2\sqrt{3}\qquad\textbf{(E) }6+\frac{\sqrt{3}}{3}$


 
 \textbf{Answer: }D



\textbf{Solution 1}

Since the triangle is equilateral and one of the sides is a vertical line, the triangle must have a horizontal line of symmetry, and therefore the other two sides will have opposite slopes. The slope of the other given line is $\frac{\sqrt{3}}{3}$ (which is must be, given 60 degree angle of the triangle, relative to vertical) so the third must be $-\frac{\sqrt{3}}{3}$. Since this third line passes through the origin, its equation is simply $y = -\frac{\sqrt{3}}{3}x$. To find two vertices of the triangle, plug in $x=1$ to both the other equations. 

 $y = -\frac{\sqrt{3}}{3}$

 $y = 1 + \frac{\sqrt{3}}{3}$

 We now have the coordinates of two vertices, $\left(1, -\frac{\sqrt{3}}{3}\right)$ and $\left(1, 1 + \frac{\sqrt{3}}{3}\right)$. The length of one side is the distance between the y-coordinates, or $1 +  \frac{2\sqrt{3}}{3}$.

 The perimeter of the triangle is thus $3\left(1 +  \frac{2\sqrt{3}}{3}\right)$, so the answer is $\boxed{\textbf{(D) }3 + 2\sqrt{3}}$

 Note: We know that the slope for the third line must be $-\frac{\sqrt{3}}{3}$, as the slope for the second given line is $\frac{\sqrt{3}}{3}$. The slope is determined from two points on a line, and we pick one of the points to be the point where they intersect. Now, drawing the symmetric line, it splits the first line; x=1 into two equal parts. Now, the second point for the third and second lines have the same x coordinate, but opposite y coordinates.

 



\textbf{Solution 1 Different Approach}

It is important to note that because \( \tan(x) \) represents the slope of a function, then \( \frac{\sqrt{3}}{3} \) is just \( \tan(30^\circ) \). Then, if we reflect \( \frac{\sqrt{3}}{3} x+1 \) over the line \( y=\frac{1}{2} \), we get \( y=-\frac{\sqrt{3}}{3}x \) as our new function.

 Given that, we may resume with solution 1.

 



\textbf{Solution 2}

Draw a line from the y-intercept of the equation  $y=1+ \frac{\sqrt{3}}{3} x$ perpendicular to the line $x=1$. There is a square of side length 1 inscribed in the equilateral triangle. The problem becomes reduced to finding the perimeter of an equilateral triangle with a square of side length 1 inscribed in it. The side length is $2\left(\frac{1}{\sqrt{3}}\right) + 1$. After multiplying the side length by 3 and rationalizing, you get $\boxed{\textbf{(D) }3 + 2\sqrt{3}}$.

 



\textbf{Solution 3}

Let the intersection point between the line $y = 1 + \frac{\sqrt{3}}{3}x$ and the line that crosses the origin be $P$.

 We drop an altitude from $P$ onto the line $x = 1$. Since the overall triangle is an equilateral triangle, we are splitting the base (on $x=1$) in half. As the y-axis is parallel to the line $x=1$, the altitude from P will also split the y-axis from $y=0$ to $y=1$ in half. From this, we can get that the y-value of P is $\frac{1}{2}$.

 Plugging this into the equation $y = 1 + \frac{\sqrt{3}}{3}x$, we get that $x = -\frac{\sqrt{3}}{2}$, and thus our height for the equilateral triangle is $1 + \frac{\sqrt{3}}{2}$. Using that, we can calculate the perimeter to be $\boxed{\textbf{(D) }3 + 2\sqrt{3}}$.

 



\textbf{Solution 4 (quick)}

$1+\frac{\sqrt{3}}{3}x$ clearly must form a 60 degree angle with $x=1$. The other line is the one that forms a $-60$ degree angle with $x=1$. Thus, it must have slope $-\frac{\sqrt{3}}{3}$, and equation $y-0=-\frac{\sqrt{3}}{3}(x-0)$.

 This equation, $x=1$, and the x-axis form a 30-60-90 triangle. The length of the short leg is $0-(-\frac{\sqrt{3}}{3}(1))$ or $\frac{\sqrt{3}}{3}$. Multiplying by 2 gives the hypotenuse and then by 3 gives the perimeter, $\boxed{\textbf{(D) }3 + 2\sqrt{3}}$

 



\textbf{Solution 5 (coord bash)}

Suppose that the line is $y=kx$. The intersections of lines $y=kx$ and $x=1$, $y=kx$ and $y=1+\frac{\sqrt{3}}{3}x$, $x=1$ and $y=1+\frac{\sqrt{3}}{3}x$, are $(1,k),\quad \left(\dfrac{3}{3k-\sqrt{3}},\dfrac{3k}{3k-\sqrt{3}}\right),\quad \left(1,\dfrac{3+\sqrt{3}}{3}\right)$ respectively. Because these 3 points form an equilateral triangle, the distances between these points are equal. Therefore, we get the equation $\left(\dfrac{3}{3k-\sqrt{3}}-1\right)^2+\left(\dfrac{3k}{3k-\sqrt{3}}-k\right)^2=\left(\dfrac{3+\sqrt{3}}{3}-k\right)^2$. Then by re-writing every term as a single fraction, we get $\left(\dfrac{3-3k+\sqrt{3}}{3k-\sqrt{3}}\right)^2+\left(\dfrac{3k-3k^2+\sqrt{3}\,k}{3k-\sqrt{3}}\right)^2=\left(\dfrac{3+\sqrt{3}-3k}{3}\right)^2$.Then we can factor out $3+\sqrt{3}-3k$ from every term, which gets us$(3+\sqrt{3}-3k)\cdot\dfrac{(3+\sqrt{3}-3k)(1+k^2)}{(3k-\sqrt{3})^2}=(3+\sqrt{3}-3k)\cdot\dfrac{3+\sqrt{3}-3k}{9}$.Cancelling out $3+\sqrt{3}-3k$ gets us $\dfrac{1+k^2}{(3k-\sqrt{3})^2}=\dfrac{1}{9}$. Solving this equation gets us

 \begin{gather*}9(1+k^2) = (3k-\sqrt{3})^2 \\[6pt]9+9k^2 = 9k^2 - 6\sqrt{3}\,k + 3 \\[6pt]9 = -6\sqrt{3}\,k + 3 \\[6pt]6 = -6\sqrt{3}\,k \\[6pt]k = -\dfrac{1}{\sqrt{3}}\end{gather*}Now we can find the perimeter of the equilateral triangle. Using the distance between $\left(1,\dfrac{3+\sqrt{3}}{3}\right)$ and $(1,k)$, the side length of the equilateral triangle is $1+\dfrac{\sqrt{3}}{3}-\left(-\dfrac{\sqrt{3}}{3}\right) = 1+\dfrac{2\sqrt{3}}{3}$. Therefore, the perimeter of the triangle is $3\left(1+\dfrac{2\sqrt{3}}{3}\right) = 3 + 2\sqrt{3}$, which gives us the answer $\boxed{\textbf{(D) }3 + 2\sqrt{3}}$

 


	\item (\textit{AMC 10 10A 2015}) Hexadecimal (base-16) numbers are written using numeric digits $0$ through $9$ as well as the letters $A$ through $F$ to represent $10$ through $15$. Among the first $1000$ positive integers, there are $n$ whose hexadecimal representation contains only numeric digits. What is the sum of the digits of $n$?


 $\textbf{(A) }17\qquad\textbf{(B) }18\qquad\textbf{(C) }19\qquad\textbf{(D) }20\qquad\textbf{(E) }21$


 
 \textbf{Answer: }E



\textbf{Solution 1}

Notice that $1000$ is $3E8$ when converted to hexadecimal ($3 \cdot 16^2 + 14 \cdot 16^1 + 8 \cdot 16^0$). We will proceed by constructing numbers that consist of only numeric digits in hexadecimal.

 The first digit could be $0,$ $1,$ $2,$ or $3,$ and the second digit and third digit could both be any digit $0 - 9$, giving $4 \cdot 10 \cdot 10 = 400$ combinations. However, this includes $000,$ so this number must be diminished by $1.$  Therefore, there are $399$ valid $n$ corresponding to those $399$ positive integers less than $1000$ that consist of only numeric digits.  (Notice that $399$ is the least hexadecimal number using only decimal digits before $3E8$.)  Therefore, our answer is $3 + 9 + 9 = \boxed{\textbf{(E) } 21}$

 



\textbf{Solution 2 (Casework)}

First, we set a bound by writing $1000$ in base-$16$. $1000_{10}=3E8_{16}$. Therefore, we are considering numbers with a maximum of $3$ digits, and a maximum of $3$ in the $256$ths-place (the first place in a $3$-digit number).

 Case $1$: $1$-digit numbers:There are evidently $9$ numbers that fit this category.

 Case $2$: $2$-digit numbers:There are $9\cdot10=90$ numbers that fit this category.

 Case $3$: $3$-digit numbers:There are $3\cdot10\cdot10=300$ numbers that fit this category

 Adding these up, we get $9+90+300=399$ numbers. $3 + 9 + 9 = \boxed{\textbf{(E) } 21}$ ~sosiaops

 



\textbf{Solution 3}

We can quickly see that $400$ in hexadecimal = $0+0+16^2*4$ = 1024. If we go down to 399 in hexadecimal, we have $9+9*16+3*256$ which is $921$, which is obviously less than 1000. Therefore, the answer is $3+9+9$ = $\boxed{\textbf{(E) } 21}$

 


	\item (\textit{AMC 10 10A 2015}) The isosceles right triangle $ABC$ has right angle at $C$ and area $12.5$. The rays trisecting $\angle ACB$ intersect $AB$ at $D$ and $E$. What is the area of $\bigtriangleup CDE$?


 $\textbf{(A) }\dfrac{5\sqrt{2}}{3}\qquad\textbf{(B) }\dfrac{50\sqrt{3}-75}{4}\qquad\textbf{(C) }\dfrac{15\sqrt{3}}{8}\qquad\textbf{(D) }\dfrac{50-25\sqrt{3}}{2}\qquad\textbf{(E) }\dfrac{25}{6}$


 
 \textbf{Answer: }D



\textbf{Solution 1 (No Trigonometry)}


                    \begin{figure}[H]
                        \centering
                        \includegraphics[width=0.5\textwidth]{Solution/AMC_10_10A_2015_19_Solution_Figure_1.png}
                    \end{figure}
                    

 $\triangle ADC$ can be split into a $45-45-90$ right triangle and a $30-60-90$ right triangle by dropping a perpendicular from $D$ to side $AC$. Let $F$ be where that perpendicular intersects $AC$.

 Because the side lengths of a $45-45-90$ right triangle are in ratio $a:a:a\sqrt{2}$, $DF = AF$.

 Because the side lengths of a $30-60-90$ right triangle are in ratio $a:a\sqrt{3}:2a$ and $AF + FC = 5$, $DF = \frac{5 - AF}{\sqrt{3}}$.

 Setting the two equations for $DF$ equal to each other, $AF = \frac{5 - AF}{\sqrt{3}}$.

 Solving gives $AF = DF = \frac{5\sqrt{3} - 5}{2}$.

 The area of $\triangle ADC =\frac12 \cdot DF \cdot AC = \frac{25\sqrt{3} - 25}{4}$.

 $\triangle ADC$ is congruent to $\triangle BEC$, so their areas are equal.

 A triangle's area can be written as the sum of the figures that make it up, so $[ABC] = [ADC] + [BEC] + [CDE]$.

 $\frac{25}{2} = \frac{25\sqrt{3} - 25}{4} + \frac{25\sqrt{3} - 25}{4} + [CDE]$.

 Solving gives $[CDE] = \frac{50 - 25\sqrt{3}}{2}$, so the answer is $\boxed{\textbf{(D) }\frac{50 - 25\sqrt{3}}{2}}$

 \textbf{Note} Another way to get $DF$ is to assume that $AF=DF$ is equal to $a$, as previously mentioned, and $CF$ is equal to $a\sqrt{3}$. $AF+DF=5=a+a\sqrt{3}$

 



\textbf{Solution 2 (Trigonometry)}

The area of $\triangle ABC$ is $12.5$, and so the leg length of $45 - 45 - 90$ $\triangle ABC$ is $5.$ Thus, the altitude to hypotenuse $AB$, $CF$, has length $\dfrac{5}{\sqrt{2}}$ by $45 - 45 - 90$ right triangles. Now, it is clear that $\angle{ACD} = \angle{BCE} = 30^\circ$, and so by the Exterior Angle Theorem, $\triangle{CDE}$ is an isosceles $30 - 75 - 75$ triangle. Thus, $DF = CF \tan 15^\circ = \dfrac{5}{\sqrt{2}} (2 - \sqrt{3})$ by the Half-Angle formula, and so the area of $\triangle CDE$ is $DF \cdot CF = \dfrac{25}{2} (2 - \sqrt{3})$. The answer is thus $\boxed{\textbf{(D) } \frac{50 - 25\sqrt{3}}{2}}$

 



\textbf{Solution 3 (Analytical Geometry)}

Because the area of triangle $ABC$ is $12.5$, and the triangle is right and isosceles, we can quickly see that the leg length of the triangle $ABC$ is 5. If we put the triangle on the coordinate plane, with vertex $C$ at the origin, and the hypotenuse in the first quadrant, we can use slope-intercept form and tangents to get three lines that intersect at the origin, $D$, and $E$. Then, you can use the distance formula to get the length of $DE$. The height is just $\frac{5}{\sqrt{2}}$, so the area is just $DE \cdot \frac{5}{\sqrt{2}} \cdot \frac{1}{2}=\boxed{\textbf{(D) } \frac{50 - 25\sqrt{3}}{2}}$

 


 



\textbf{Solution 4 (Weak Trigonometry)}

Just like with Solution 1, we drop a perpendicular from $D$ onto $AC$, splitting it into a $30$-$60$-$90$ triangle and a $45$-$45$-$90$ triangle. We find the legs of the $45$-$45$-$90$ triangle are $x$ and $x$. The other leg of the other triangle is $5-x$. By  using $x$ as the short leg of $30$-$60$-$90$ triangle, we find that the long leg is equal to $x\sqrt{3}$. We then set $5-x$ = $x$ $\sqrt{3}$. So, $5 = x(\sqrt{3}+1) x = \frac{5}{(\sqrt{3}+1)}$ Rationalize by multiplying both numerator and denominator by $\sqrt{3}-1$  we get $5(\frac{(\sqrt{3}-1)}{(\sqrt{3})^2-1^2})$. Let $x$ be AF. We find that $AF=\frac{5\sqrt{3}-5}{2}$.

 Now, since $\triangle AEC\cong \triangle BDC$ by ASA, $CE=CD$. Since, $DF=\frac{5\sqrt{3}-5}{2}$, $DC=2\cdot \frac{5\sqrt{3}-5}{2}=5\sqrt{3}-5$. By the sine area formula, $[CDE]=\frac{1}{2}\cdot \sin 30\cdot CD^2=\frac{1}{4}\cdot (100-50\sqrt{3})=\frac{50-25\sqrt{3}}{2}\implies \boxed{\textbf{(D)}}$

 



\textbf{Solution 5 (Basic Trigonometry)}

Prerequisite knowledge for this solution: the side ratios of a 30-60-90, and 45-45-90 right triangle.

 
We let point C be the origin. Since $\overline{CD}$ and $\overline{CE}$ trisect $\angle ACB = 90^{\circ}$, this means $m\angle CEB = 30^{\circ}$ and the equation of $\overline{CE}$ is $y=\frac{\sqrt{3}}{3}$ (you can figure out that the tangent of 30 degrees gives $\frac{\sqrt{3}}{3}$). Next, we can find A to be at $(0, 5)$ and B at $(5, 0)$, so the equation of $\overline{AB}$ is $y=-x+5$. So we have the system:

 
 \[\begin{cases}y=\frac{\sqrt{3}}{3}x\\y=-x+5\end{cases}\]

 By substituting values, we can arrive at $\frac{3+\sqrt{3}}{3}x=5$, or $x=5\cdot\frac{3}{3+\sqrt{3}}=\frac{15}{3+\sqrt{3}}$. We multiply $x=\frac{15}{3+\sqrt{3}}\cdot\frac{3-\sqrt{3}}{3-\sqrt{3}}=\frac{45-15\sqrt{3}}{6}=\frac{15-5\sqrt{3}}{2}$.

 Dropping an altitude from E onto $\overline{CB}$, and calling the intersection point G, we find that $\triangle EGB$ is a 45-45-90 triangle with a leg of $\frac{15-5\sqrt{3}}{2}\cdot\frac{\sqrt{3}}{3}=\frac{15\sqrt{3}-15}{6}=\frac{5\sqrt{3}-5}{2}$. Thus, $EB=\frac{5\sqrt{3}-5}{2}\sqrt{2}=\frac{5\sqrt{6}-5\sqrt{2}}{2}$.

 Dropping an altitude from C onto $\overline{AB}$, and calling the intersection point H, we find that $CH=\frac{5\sqrt{2}}{2}=BH$, and by the Theorem of Betweenness applied to H, E, and B, we get $HE=HB-EB=\frac{5\sqrt{2}}{2}-\frac{5\sqrt{6}-5\sqrt{2}}{2}=\frac{10\sqrt{2}-5\sqrt{6}}{2}$.

 We are almost done. By symmetry, $HD=HE$, so to find the area of the triangle CED, we only need to multiply HE by CH, $\frac{10\sqrt{2}-5\sqrt{6}}{2}\cdot\frac{5\sqrt{2}}{2}=\frac{100-50\sqrt{3}}{4}=\frac{50-25\sqrt{3}}{2}$. This is answer choice $\boxed{\textbf{(D) } \frac{50-25\sqrt{3}}{2}}$~JH. L~Minor edit by Mathcrafter314

 



\textbf{Solution 6 (Ratio of 15-75-90 Triangle)}

Let $\triangle ABC$ be an isosceles right triangle with the right angle at vertex $C$. Two rays from $C$ trisect $\angle C$ into three $30^\circ$ angles, intersecting the hypotenuse $AB$ at $D$ and $E$. 

 By symmetry, $\triangle CAD \cong \triangle CBE$. Because $\angle A = 45^\circ$ and $\angle ACD = 30^\circ$, the exterior angle theorem gives $\angle CDE = 45^\circ + 30^\circ = 75^\circ$. Similarly, $\angle CED = 75^\circ$, making $\triangle CDE$ an isosceles triangle with a vertex angle of $30^\circ$.

 Let the equal legs of the isosceles right triangle be $AC = BC = s$. The area of $\triangle ABC$ is given as $12.5$, so:
 \[\text{Area}(\triangle ABC) = \frac{1}{2}s^2 = 12.5 \implies s^2 = 25 \implies s = 5\]

 By the Pythagorean theorem, the hypotenuse $AB$ is:
 \[AB = \sqrt{s^2 + s^2} = s\sqrt{2} = 5\sqrt{2}\]

 Let $CH$ be the altitude from $C$ to $AB$, and let its length be $x$. Because $\triangle ABC$ is an isosceles right triangle, the altitude splits the hypotenuse into two equal halves ($AH = BH = \frac{5\sqrt{2}}{2}$) and creates a $45^\circ$-$45^\circ$-$90^\circ$ right triangle $\triangle ACH$. Thus, the altitude is equal to half the hypotenuse:
 \[x = CH = \frac{5\sqrt{2}}{2}\]

 The altitude $CH$ bisects the central $30^\circ$ angle $\angle DCE$, splitting $\triangle CDE$ into two identical $15^\circ$-$75^\circ$-$90^\circ$ right triangles ($\triangle CHD$ and $\triangle CHE$) where $\angle DCH = \angle ECH = 15^\circ$. 

 From the known side-length ratio of a $15^\circ$-$75^\circ$-$90^\circ$ triangle, the ratio of the shorter leg (opposite the $15^\circ$ angle) to the longer leg (opposite the $75^\circ$ angle) is exactly $2 - \sqrt{3}$. Therefore, the bottom leg $DH$ is:
 \[DH = CH \times (2 - \sqrt{3}) = x(2 - \sqrt{3})\]

 Since $H$ is the midpoint of $DE$, the full base of the central triangle is $DE = 2DH = 2x(2 - \sqrt{3})$. The area of $\triangle CDE$ can be computed by doubling the area of right triangle $\triangle CHE$:
 \[\text{Area}(\triangle CDE) = 2 \times \text{Area}(\triangle CHE) = 2 \times \left( \frac{1}{2} \cdot DH \cdot CH \right) = DH \cdot CH\]

 Substituting our expressions into the area formula yields:
 \[\text{Area}(\triangle CDE) = [x(2 - \sqrt{3})] \cdot x = x^2(2 - \sqrt{3})\]

 Since $x = \frac{5\sqrt{2}}{2}$, we find $x^2 = \frac{25 \cdot 2}{4} = \frac{25}{2}$. Substituting this value gives the final exact area:
 \[\text{Area}(\triangle CDE) = \frac{25}{2}(2 - \sqrt{3}) = \frac{50-25\sqrt{3}}{2}\]

 Therefore, the final answer is $\boxed{\textbf{(D) } \frac{50-25\sqrt{3}}{2}}$

 


 \textbf{Diagram} Below is the geometric visualization displaying all calculated angle measures:

 
                    \begin{figure}[H]
                        \centering
                        \includegraphics[width=0.5\textwidth]{Solution/AMC_10_10A_2015_19_Solution_Figure_2.png}
                    \end{figure}
                    

 



\textbf{Solution 7 (Law of Sines)}

We know that the area of the right triangle is $\frac{25}{2}$ and that the two legs are equal, so we can easily tell that the length of the two legs is $5$. Thus, the hypotenuse $AB = 5\sqrt{2}$ and $\angle{CAB} = \angle{ABC} = 45^{\circ}.$Let's quickly define $F$ as the point that bisects $\angle{ACB}$ and $\overline{AB}$. Then, we can say that the area of the desired triangle is $DF \cdot AF$.Let $\overline{AD} = \overline{BE} = n.$ Since $D$ is one of the trisecting points of $\angle{ACB}, \angle{ACD} = 30^{\circ}.$ Because 
 \[\angle{ADC} = 180^{\circ}-\angle{ACD}-\angle{CAB},\]
 \[\angle{ADC} = 180^{\circ} - 45^{\circ} - 30^{\circ} = 105^{\circ}.\]Now, we can employ the Law of Sines. It tells us that $\frac{\sin(\angle{ADC})}{5} = \frac{\sin(\angle{ACD})}{n}$. Plugging in our angle values, we get that 
 \[\frac{\sin(105^{\circ})}{5} = \frac{\sin(30^{\circ})}{n}.\] It's easy to find that $\sin(105^{\circ}) = \sin(75^{\circ}) = \frac{\sqrt{2}+\sqrt{6}}{4},$ and that $\sin(30^{\circ}) = \frac{1}{2}$. Plugging in these values into our previous equation, we get 
 \[\frac{\sqrt{2}+\sqrt{6}}{20} = \frac{1}{2n}.\] Cross multiplying gets us 
 \[2n\sqrt{2}+2n\sqrt{6} = 20,\] and then we simplify like so: 
 \[2n(\sqrt{2}+\sqrt{6}) = 20\rightarrow\]
 \[n(\sqrt{2}+\sqrt{6}) = 10\rightarrow\]
 \[n = \frac{10}{\sqrt{2}+\sqrt{6}}.\]Now, using our definition of $n$, we know that $DF$ = $\frac{AB}{2} - n = \frac{5\sqrt{2}}{2} - \frac{10}{\sqrt{2}+\sqrt{6}}$. We want to put this under one common denominator, which is pretty simple to execute. That leaves us with 
 \[\frac{5\sqrt{2} \cdot \sqrt{2}+5\sqrt{2} \cdot \sqrt{6} - 20}{2\sqrt{2} + 2\sqrt{6}}=\]
 \[\frac{10+10\sqrt{3} - 20}{2\sqrt{2} + 2\sqrt{6}}=\]
 \[\frac{10\sqrt{3} - 10}{2(\sqrt{2} + \sqrt{6})}=\]
 \[\frac{2(5\sqrt{3} - 5)}{2(\sqrt{2} + \sqrt{6})}=\]
 \[\frac{5\sqrt{3} - 5}{\sqrt{2} + \sqrt{6}}.\]Whew. That was longer than expected. Anyways, quick inspection tells us that $AF = \frac{5\sqrt{2}}{2},$ so now we just have to do some simplifying to find the desired, $[CDE]$. Let's do that now.
 \[[CDE] = \frac{5\sqrt{3} - 5}{\sqrt{2} + \sqrt{6}} \cdot \frac{5\sqrt{2}}{2}=\]
 \[[CDE] = \frac{25\sqrt{6} - 25\sqrt{2}}{2(\sqrt{2} + \sqrt{6})}=\]
 \[[CDE] = \frac{25(\sqrt{2} - \sqrt{6})}{2(\sqrt{2} + \sqrt{6})}=\](We need to take a quick conjugation break. Note that $(\sqrt{2} - \sqrt{6})^2 = 8 - 4\sqrt{3}.$)
 \[[CDE] = \frac{25(\sqrt{2} - \sqrt{6})}{2(\sqrt{2} + \sqrt{6})} \cdot \frac{(\sqrt{2} - \sqrt{6})}{(\sqrt{2} - \sqrt{6})}=\]
 \[[CDE] = \frac{200 - 100\sqrt{3}}{2(4)}=\]
 \[[CDE] =\boxed{\frac{50 - 25\sqrt{3}}{2} \text{(D)}.}\]~Nickelslordm

 


 \textbf{Side Ratio Proof Diagram} This auxiliary construction extends a standard $30^\circ$-$60^\circ$-$90^\circ$ right triangle ($\triangle XYZ$) to create a larger $15^\circ$-$75^\circ$-$90^\circ$ right triangle ($\triangle XYW$). This layout allows calculation of the side ratios without using trigonometry functions:

 
                    \begin{figure}[H]
                        \centering
                        \includegraphics[width=0.5\textwidth]{Solution/AMC_10_10A_2015_19_Solution_Figure_3.png}
                    \end{figure}
                    

 


	\item (\textit{AMC 10 10A 2015}) A rectangle with positive integer side lengths in $\text{cm}$ has area $A$ $\text{cm}^2$ and perimeter $P$ $\text{cm}$. Which of the following numbers cannot equal $A+P$?


 $\textbf{(A) }100\qquad\textbf{(B) }102\qquad\textbf{(C) }104\qquad\textbf{(D) }106\qquad\textbf{(E) }108$


 \textbf{Note:} As it originally appeared in the AMC 10, this problem was stated incorrectly and had no answer; it has been modified here to be solvable. This is the original question:


 A rectangle with side lengths in $\text{cm}$ has an area of integer $A$ $\text{cm}^2$ and a perimeter of integer $P$ $\text{cm}$. Which of the following numbers cannot equal $A+P$?


 $\textbf{(A) }100\qquad\textbf{(B) }102\qquad\textbf{(C) }104\qquad\textbf{(D) }106\qquad\textbf{(E) }108$


 
 \textbf{Answer: }B (Note: This problem was originally stated incorrectly, and all contestants received full credit regardless of their answer.)



\textbf{Solution 1 (Simon's Favorite Factoring Trick)}

Let the rectangle's length be $a$ and its width be $b$. Its area is $ab$ and the perimeter is $2a+2b$.

 Then $A + P = ab + 2a + 2b$. Factoring, we have $(a + 2)(b + 2) - 4$.

 The only one of the answer choices that cannot be expressed in this form is $102$, as $102 + 4$ is twice a prime. There would then be no way to express $106$ as $(a + 2)(b + 2)$, keeping $a$ and $b$ as positive integers.

 Our answer is then $\boxed{B}$.

 Note: The original problem only stated that $A$ and $P$ were positive integers, not the side lengths themselves. This rendered the problem unsolvable, and so the AMC awarded everyone 6 points on this problem. This wiki has the corrected version of the problem so that the 2015 AMC 10A test can be used for practice.

 


	\item (\textit{AMC 10 10A 2015}) Tetrahedron $ABCD$ has $AB=5$, $AC=3$, $BC=4$, $BD=4$, $AD=3$, and $CD=\frac{12}5\sqrt2$.  What is the volume of the tetrahedron?


 $\textbf{(A) }3\sqrt2\qquad\textbf{(B) }2\sqrt5\qquad\textbf{(C) }\dfrac{24}5\qquad\textbf{(D) }3\sqrt3\qquad\textbf{(E) }\dfrac{24}5\sqrt2$


 
 \textbf{Answer: }C



\textbf{Solution 1}

Drop altitudes of triangle $ABC$ and triangle $ABD$ down from $C$ and $D$, respectively. Both will hit the same point; let this point be $T$. Because both triangle $ABC$ and triangle $ABD$ are 3-4-5 triangles, $CT = DT = \dfrac{3\cdot4}{5} = \dfrac{12}{5}$. Because $CT^{2} + DT^{2} = 2\left(\frac{12}{5}\right)^{2} = \left(\frac{12}{5}\sqrt{2}\right)^{2} = CD^{2}$, it follows that the $CTD$ is a right triangle, meaning that $\angle CTD = 90^\circ$, and it follows that planes $ABC$ and $ABD$ are perpendicular to each other. Now, we can treat $ABC$ as the base of the tetrahedron and $TD$ as the height. Thus, the desired volume is 
 \[V = \dfrac{1}{3} bh = \dfrac{1}{3}\cdot[ABC]\cdot TD = \dfrac{1}{3} \cdot 6 \cdot \dfrac{12}{5} = \dfrac{24}{5}\] which is answer $\boxed{\textbf{(C) } \dfrac{24}{5}}$

 



\textbf{Solution 2}

Let the midpoint of $CD$ be $E$. We have $CE = \dfrac{6}{5} \sqrt{2}$, and so by the Pythagorean Theorem $AE = \dfrac{\sqrt{153}}{5}$ and $BE = \dfrac{\sqrt{328}}{5}$. Because the altitude from $A$ of tetrahedron $ABCD$ passes touches plane $BCD$ on $BE$, it is also an altitude of triangle $ABE$. The area $A$ of triangle $ABE$ is, by Heron's Formula, given by

 
 \[16A^2 = 2a^2 b^2 + 2b^2 c^2 + 2c^2 a^2 - a^4 - b^4 - c^4 = -(a^2 + b^2 - c^2)^2 + 4a^2 b^2.\]Substituting $a = AE, b = BE, c = 5$ and performing huge (but manageable) computations yield $A^2 = 18$, so $A = 3\sqrt{2}$. Thus, if $h$ is the length of the altitude from $A$ of the tetrahedron, $BE \cdot h = 2A = 6\sqrt{2}$. Our answer is thus
 \[V = \dfrac{1}{3} Bh = \dfrac{1}{3} h \cdot BE \cdot \dfrac{6\sqrt{2}}{5} = \dfrac{24}{5},\]and so our answer is $\boxed{\textbf{(C) } \dfrac{24}{5}}$

 



\textbf{Solution 3}

Similar to solution 1, $\triangle CTD$ is an isosceles right triangle. $AB$ is perpendicular to the plane $CTD$. So, we can cut tetrahedron $ABCD$ into 2 tetrahedrons $ACTD$ and $BCTD$ with $\triangle CTD$ as their common base, $BT$ and $AT$ as their heights.

 
$[CTD]=\frac{1}{2} \cdot CT \cdot DT=\frac{1}{2} \cdot \frac{12}{5} \cdot \frac{12}{5}= \frac{72}{25}$

 $V_{ABCD}=V_{ACDT}+V_{BCDT}=\frac{1}{3} \left([CTD] \cdot AT + [CTD] \cdot BT \right)=\frac{1}{3} \cdot [CTD] \cdot (AT+BT)=\frac{1}{3} \cdot [CTD] \cdot AB  =\frac{1}{3} \cdot \frac{72}{25} \cdot 5$

 $= \boxed{\textbf{(C) } \dfrac{24}{5}}$

 


 



\textbf{Solution 4 (Francesca's Irregular Tetrahedron Formula)}

Note: Please don't ever try doing this in an actual competition. It's fun to do, however.

 Using Piero della Francesca's Theorem:
 \[V = \frac{1}{12} \left( \sqrt{-a^2b^2c^2 - a^2d^2e^2 - b^2d^2f^2 - c^2e^2f^2 + a^2c^2d^2 + b^2c^2d^2 + a^2b^2e^2 + b^2c^2e^2 + b^2d^2e^2 + c^2d^2e^2 + a^2b^2f^2 + a^2c^2f^2 + a^2d^2f^2 + c^2d^2f^2 + a^2e^2f^2 + b^2e^2f^2 - c^4d^2 - c^2d^4 - b^4e^2 - b^2e^4 - a^4f^2 - a^2f^4} \right)\]Substituting this all in... You get$\boxed{\textbf{(C) } \dfrac{24}{5}}$

 



\textbf{Solution 5 (Coord Bash)}

Set $B$ as $(0,0,0)$, $C$ at $(4,0,0)$, and $A$ at $(4,3,0)$. If we find $D$'s coordinates then the z-coordinate(the last coord) is the height we need since the rest of the coordinates have a z-coord of $0$. Notice that we are given every distance from point $D$, we can use basically satellite-type triangulation. $D$ is at every point a radius from points $A,B$ and $C$, or a sphere with the center at points $A,B$ and $C$. We can set up a system of equations to find where all 3 spheres intersect. 
 \[\begin{split} (x-4)^2+(y-3)^2+z^2 &=3^2 \\ x^2+y^2+z^2 &=4^2 \\ (x-4)^2+y^2+z^2 &=(\frac{12}{5}\sqrt2)^2 \end{split}\]Solving this system of equations gives you:
 \[\begin{split} x &=\frac{64}{25} \\ y &=\frac{48}{25} \\ z &=\frac{12}{5}  \end{split}\]Since $z$ is the value we need, we can now solve for the volume. The base is just $\frac{3*4}{2}=6$
 \[\begin{split} V &=\frac{1}{3}bh \\ V &=\frac{1}{3}(6)(\frac{12}{5}) \\ V &=\frac{72}{15} \\ V &=\boxed{\textbf{(C)}\dfrac{24}{5}} \end{split}\]

 


	\item (\textit{AMC 10 10A 2015}) Eight people are sitting around a circular table, each holding a fair coin. All eight people flip their coins and those who flip heads stand while those who flip tails remain seated. What is the probability that no two adjacent people will stand?


 $\textbf{(A)}\dfrac{47}{256}\qquad\textbf{(B)}\dfrac{3}{16}\qquad\textbf{(C) }\dfrac{49}{256}\qquad\textbf{(D) }\dfrac{25}{128}\qquad\textbf{(E) }\dfrac{51}{256}$


 
 \textbf{Answer: }A



\textbf{Solution 1}

We will count how many valid standing arrangements there are (counting rotations as distinct), and divide by $2^8 = 256$ at the end. We casework on how many people are standing.

 Case $1:$ $0$ people are standing. This yields $1$ arrangement.

 Case $2:$ $1$ person is standing. This yields $8$ arrangements.

 Case $3:$ $2$ people are standing. This yields $\dbinom{8}{2} - 8 = 20$ arrangements, because the two people cannot be next to each other.

 Case $5:$ $4$ people are standing. Then the people must be arranged in stand-sit-stand-sit-stand-sit-stand-sit fashion, yielding $2$ possible arrangements.

 More difficult is:

 Case $4:$ $3$ people are standing. First, choose the location of the first person standing ($8$ choices). Next, choose $2$ of the remaining people in the remaining $5$ legal seats to stand, amounting to $6$ arrangements considering that these two people cannot stand next to each other. However, we have to divide by $3,$ because there are $3$ ways to choose the first person given any three. This yields $\dfrac{8 \cdot 6}{3} = 16$ arrangements for Case $4.$

 Alternate Case $4:$ Use complementary counting. Total number of ways to choose 3 people from 8 which is $\dbinom{8}{3}$. Sub-case $1:$ three people are next to each other which is $\dbinom{8}{1}$. Sub-case $2:$ two people are next to each other and the third person is not $\dbinom{8}{1}$ $\dbinom{4}{1}$. This yields $\dbinom{8}{3} - \dbinom{8}{1} - \dbinom{8}{1} \dbinom{4}{1} = 16$ 

 Summing gives $1 + 8 + 20 + 2 + 16 = 47,$ and so our probability is $\boxed{\textbf{(A) } \dfrac{47}{256}}$.

 Alternate: (Quicksolve) - We know that for case 5, there are 8 ways to choose the first person, and 3 ways to choose the first person given any 3 - which means that for the case, there is 8 * something divided by 3. The sum of the other cases is 31/256. Thus, add a multiple of 8 to 31 to get an answer. The options are 31 + 8 = 39, 31 + 16 = 47, 31 + 24 = 55, etc. The only possible answer is 47/256.

 



\textbf{Solution 1.5 (solution 1 but with visuals)}

\textbf{Case 1:} Everyone flips tails

 1 way:
                    \begin{figure}[H]
                        \centering
                        \includegraphics[width=0.5\textwidth]{Solution/AMC_10_10A_2015_22_Solution_Figure_1.png}
                    \end{figure}
                    

 
\textbf{Case 2:} One person flips heads 

 8 ways:
                    \begin{figure}[H]
                        \centering
                        \includegraphics[width=0.5\textwidth]{Solution/AMC_10_10A_2015_22_Solution_Figure_2.png}
                    \end{figure}
                    

 \textbf{Case 3:} Two people flip heads 

 There are $\binom{8}{2}$ ways to choose any $2$ people in the circle and $8$ ways to choose $2$ people adjacent to each other. So this case gives us $\binom{8}{2} - 8 = 20$ ways. An example of choosing two nonadjacent people is shown below. 

 
                    \begin{figure}[H]
                        \centering
                        \includegraphics[width=0.5\textwidth]{Solution/AMC_10_10A_2015_22_Solution_Figure_3.png}
                    \end{figure}
                    

 
\textbf{Case 4:} Three people flip heads

 We can choose the first person (shaded gray) in $8$ ways. We can choose the other $2$ people in $6$ ways (purple-blue, yellow-red, green-yellow, green-blue, purple-red, green-red). We divide by $3$ for overcounting to get $\frac{8 \cdot 6}{3} = 16$ ways for this case. 

 
                    \begin{figure}[H]
                        \centering
                        \includegraphics[width=0.5\textwidth]{Solution/AMC_10_10A_2015_22_Solution_Figure_4.png}
                    \end{figure}
                    

 \textbf{Case 5:} 4 people flip heads

 2 ways (we can either choose the people shaded in red or blue)

 
                    \begin{figure}[H]
                        \centering
                        \includegraphics[width=0.5\textwidth]{Solution/AMC_10_10A_2015_22_Solution_Figure_5.png}
                    \end{figure}
                    

 It's impossible to have more than 4 people standing without anyone being next to each other, so we stop here. 

 The total number of ways is $2^8,$ so the answer is $\frac{1 + 8 + 20 + 16 + 2}{2^8} = \boxed{\frac{47}{256}}.$

 \textbf{Useful Note (DARSHSHETTY4)} It is usually helpful to fix one person (element) as the point of view to simplify such circular arrangement problems.

 



\textbf{Solution 2}

We will count how many valid standing arrangements there are counting rotations as distinct and divide by $256$ at the end.Line up all $8$ people linearly.  In order for no two people standing to be adjacent, we will place a sitting person to the right of each standing person.  In effect, each standing person requires $2$ spaces and the standing people are separated by sitting people.  We just need to determine the number of combinations of pairs and singles and the problem becomes very similar to pirates and gold aka stars and bars aka sticks and stones aka balls and urns.

 If there are $4$ standing, there are ${4 \choose 4}=1$ ways to place them.For $3,$ there are ${3+2 \choose 3}=10$ ways.etc.Summing, we get ${4 \choose 4}+{5 \choose 3}+{6 \choose 2}+{7 \choose 1}+{8 \choose 0}=1+10+15+7+1=34$ ways.

 Now we consider that the far right person can be standing as well, so we have${3 \choose 3}+{4 \choose 2}+{5 \choose 1}+{6 \choose 0}=1+6+5+1=13$ ways

 Together we have $34+13=47$, and so our probability is $\boxed{\textbf{(A) } \dfrac{47}{256}}$.

 



\textbf{Solution 3 (Recursion)}

Because the denominator is always $256$, we can count the numerator with a recursion. Define $c_n$ and $r_n$ as the number of satisfying arrangements on a circle and row with $n$ seats, respectively. 

 Then, observe that 
 \[c_n=r_{n-1}+r_{n-3}\] 
 \[r_n=r_{n-1}+r_{n-2}\](From casework on whether the $n$th position has $T$ or $H$.). 

 Because $r_0=1, r_1=2$, we can continue to write out each $r_i$ and hence find $c_8=47$. Our answer is then $\boxed{\textbf{(A) } \dfrac{47}{256}}$.

 



\textbf{Solution 4 (Recursion)}

We know that the denominator of the probability is $2^8=256$. So now we only have to calculate the numerator, which is the number of arrangements for $8$ people at a round table without $2$ or more neighboring people standing.

 Denote $a_n$ as number of arrangements for $n$ people at a round table without $2$ or more neighboring people standing. We can see that $a_2=3$, $a_3=4$, we are going to prove $a_n=a_{n-1}+a_{n-2}$

 We use $1$ to represent standing people, and $0$ as sitting people. The problem becomes arranging $n$ numbers of $0$ or $1$ around a circle with no consecutive $1$'s.

 We elect one person as the chairman, let him be $p_1$, the first element in this circular sequence.There are $2$ cases to generate the arrangement of $n$ numbers.

 $Case$ $1:$ From the arrangement of $n-1$ numbers, add $1$ more number $p_n$ counter-clockwise next to $p_1$. 

 If $p_1=1$, $p_{n-1}p_np_1$ is $001$, as the diagram shows:

 
                    \begin{figure}[H]
                        \centering
                        \includegraphics[width=0.5\textwidth]{Solution/AMC_10_10A_2015_22_Solution_Figure_6.png}
                    \end{figure}
                    

 If $p_1=0$, $p_{n-1}p_np_1$ is $X00$, $X$ could be $0$ or $1$, as the diagram shows:

 
                    \begin{figure}[H]
                        \centering
                        \includegraphics[width=0.5\textwidth]{Solution/AMC_10_10A_2015_22_Solution_Figure_7.png}
                    \end{figure}
                    

 $Case$ $2:$ From the arrangement of $n-2$ numbers, add $2$ more numbers $p_{n-1}$ and $p_n$ counter-clockwise next to $p_1$.

 If $p_1=1$, then $p_{n-2}=0$, let $p_{n-1}=1, p_n=0$, $p_{n-1}p_np_1$ is $101$, as the diagram shows:

 
                    \begin{figure}[H]
                        \centering
                        \includegraphics[width=0.5\textwidth]{Solution/AMC_10_10A_2015_22_Solution_Figure_8.png}
                    \end{figure}
                    

 If $p_1=0$, then $p_{n-2}$ could be $0$ or $1$. Let $p_{n-1}=0, p_n=1$, $p_{n-1}p_np_1$ is $010$, as the diagram shows:

 
                    \begin{figure}[H]
                        \centering
                        \includegraphics[width=0.5\textwidth]{Solution/AMC_10_10A_2015_22_Solution_Figure_9.png}
                    \end{figure}
                    

 From the above $2$ cases and the $4$ diagrams, the arrangements of $p_{n-1}p_np_1$ are mutually exclusive collectively exhaustive, so $a_{n}=a_{n-1}+a_{n-2}$

 \begin{verbatim}
$a_2=3$$a_3=4$$a_4=7$$a_5=11$$a_6=18$$a_7=29$$a_8=47$
\end{verbatim}

The answer is $\boxed{\textbf{(A) } \dfrac{47}{256}}$

 This sequence of numbers is called \href{https://en.wikipedia.org/wiki/Lucas\_number}{Lucas Number}.

 



\textbf{Solution 5 (Recursion)}

Note that there are $2^8=256$ cases in total, so it suffices to find the number of satisfying cases. Define $h_n$ and $t_n$ as sequences of $h$’s and $t$’s where no $2$ $h$’s are adjacent. The sequences will end in $h$ and $t$, respectively (ignore the fact that they are in a circle for now).

 In order for a sequence to end in $h$, the previous term has to be a $t$; a sequence with a length of $n$ ending in $h$ is equivalent to adding $h$ to the end of a sequence of length $n-1$ that ends with $t$. This means $h_n=t_{n-1}$. 

 In order for a sequence to end in $t$, the previous term can be $h$ or $t$, so $t_n=t_{n-1}+h_{n-1}$.

 Now, we have to take account of the circle. Notice that $t_n$ does not depend on if the people are in a circle or not, but $h_n$ does. To guarantee that some sequence that ends with $h$ does not start with $h$, we have to add a $t$ in the beginning of the sequence. Since there are $n-1$ terms left, there are $h_{n-1}$ sequences which end with $h$ that satisfy the circle condition.

 The length of our sequence is $8$, so we need to find $t_8+h_7$. Manipulating our recursions, we can find that $t_n=t_{n-1}+t_{n-2}$. Now, we proceed to find our terms.

 $t_1=1$

 $t_2=2$

 $t_3=3$

 $t_4=5$

 $t_5=8$

 $t_6=13$

 $t_7=21$

 $t_8=34$

 Since $h_7=t_6$ due to the fact that $h_n=t_{n-1}$, our answer is $\frac{34+13}{256}=\boxed{\textbf{(A) } \dfrac{47}{256}}$.

 



\textbf{Solution 6 (big fas)}

There can be 0, 1, 2, 3, or 4 people standing up. 

 0 people standing up : 1 way

 
1 person standing up : 8 ways

 
2 people standing up : 8 for the first, 5 for the second, divide by 2 for double count, 20 ways.

 
3 people standing up : we can start drawing, and notice we can only have gaps of (1,1,3) or (2,2,1). Rotating, we get 8+8=16.

 
4 people standing up : must be alternating all the way through. 2 ways

 $1+8+20+16+2 = 47$, since there are $2^8=256$ ways total, the answer is $\boxed{\textbf{(A)}\frac{47}{256}}$

 -skibbysiggy

 \textbf{Soluion 7} A nice trick when it comes to counting problems involving circles with people surrounded around the circumference of the circle is to map the circle onto a row. In other words, the problem is reduced to the below exact same problem:

 Consider a row of $8$ people in which each person is either named $H$ or $T$. If the first and last people both cannot be named $H$ at the same time and no two consecutive people are named $H$, how many possible scenarios are there?

 We have two possible starting cases:

 H T \_ \_ \_ \_ \_ T

 T \_ \_ \_ \_ \_ \_ \_

 The first case is like this since after the first $H$ we need to put a $T$ and the last can't be the same as the first, which is $H$, so it must be $T$. The second case has no restrictions yet. 

 Let $f(n)$ be the number of ways for $n$ people seated in a row. We're looking for $f(8)$. Let $g(n)$ be the number of ways with the same restrictions as $f(n)$ except that there is no restriction that the first and last person can't have both names of $H$ at the same time. Therefore, we write

 $f(8) = g(5) + g(7)$

 Consider $n$ people seated in a row. We want to find a recursive definition for $g(n)$. We again consider two cases:

 T \_ \_ ... \_ \_

 H T \_ \_ ... \_ \_

 Thus, $g(n) = g(n - 1) + g(n - 2)$ and so we can write

 $g(5) + g(7)$ in a Fibonacci type way.

 Consider $g(1)$. There's no restrictions as there's only one person which could be named either $H$ or $T$ so $g(1) = 2$.$g(2) = 3$ because we only have

 $HT, TH, TT$

 Now we use our recursive definition to find $g(5) + g(7)$. We have

 $g(3) = g(2) + g(1) = 3 + 2 = 5$

 $g(4) = g(3) + g(2) = 5 + 3 = 8$

 $g(5) = g(4) + g(3) = 8 + 5 = 13$

 $g(6) = g(5) + g(4) = 13 + 8 = 21$

 $g(7) = g(6) + g(5) = 21 + 13 = 34$

 The answer is $f(8) = g(5) + g(7) = 13 + 34 = 47$. 

 But going back to the original problem, we're looking for the probability. There's a total of $2^8 = 256$ possibilities since each of the $8$ people can be named either $H$ or $T$. The final answer is 

 $\frac{47}{256}$

 


 



\textbf{Solution 8 (Complementary Counting via PIE)}

We will find the number of invalid configurations (where at least one pair of adjacent people are both standing) and subtract this from the total possible coin-flip outcomes, which is $2^8 = 256$.

 Let the 8 seats around the table be numbered $1, 2, \dots, 8$ in order. A configuration is considered "bad" if any of the 8 adjacent pairs—$(1,2), (2,3), \dots, (7,8),$ or $(8,1)$—consist of two standing people. Let $A_i$ represent the set of configurations where the adjacent pair starting at seat $i$ are both standing. The total number of bad ways is the size of the union of these sets:
 \[|A_1 \cup A_2 \cup \dots \cup A_8|\]

 We calculate this union using the Principle of Inclusion-Exclusion (PIE):

 \textbf{Step 1: Single adjacent pairs ($\sum |A_i|$)}There are $\dbinom{8}{1} = 8$ choices for a single bad pair. Forcing these 2 specific people to stand leaves the other 6 people free to flip either heads or tails.
 \[8 \times 2^6 = 8 \times 64 = 512\]

 \textbf{Step 2: Intersections of two pairs ($\sum |A_i \cap A_j|$)}There are $\dbinom{8}{2} = 28$ ways to select two pairs. We categorize them by whether they share a person:

 
                \begin{itemize}
                    \item Overlapping pairs: The pairs share a person (e.g., $1\text{-}2$ and $2\text{-}3$). This forces 3 consecutive people to stand, leaving $8 - 3 = 5$ people free. There are $8$ such overlapping pairs.
                \end{itemize}


                
 \[8 \times 2^5 = 256\]

 
                \begin{itemize}
                    \item Disjoint pairs: The two pairs do not share any people (e.g., $1\text{-}2$ and $4\text{-}5$). This forces 4 distinct people to stand, leaving $8 - 4 = 4$ people free. There are $28 - 8 = 20$ such pairs.
                \end{itemize}


                
 \[20 \times 2^4 = 320\]Summing these gives $256 + 320 = 576$.

 \textbf{Step 3: Intersections of three pairs ($\sum |A_i \cap A_j \cap A_k|$)}There are $\dbinom{8}{3} = 56$ ways to select three pairs. We organize them by how many distinct people are forced to stand:

 
                \begin{itemize}
                    \item 4 people forced to stand: A single continuous chain of 3 bad pairs (e.g., $1\text{-}2, 2\text{-}3, 3\text{-}4$). There are $8$ ways, leaving 4 people free: $8 \times 2^4 = 128$.\item 5 people forced to stand: One overlapping block of 2 pairs and one separate isolated pair (e.g., $1\text{-}2\text{-}3$ and $5\text{-}6$). There are $8 \times 4 = 32$ ways, leaving 3 people free: $32 \times 2^3 = 256$.\item 6 people forced to stand: Three completely separated, disjoint pairs (e.g., $1\text{-}2$, $4\text{-}5$, $7\text{-}8$). There are $16$ ways to rotate and shift this setup, leaving 2 people free: $16 \times 2^2 = 64$.
                \end{itemize}


                Summing these gives $128 + 256 + 64 = 448$.

 \textbf{Step 4: Intersections of four pairs ($\sum |A_i \cap \dots \cap A_l|$)}There are $\dbinom{8}{4} = 70$ total combinations of four pairs:

 
                \begin{itemize}
                    \item Continuous chain of 4 pairs (forces 5 people to stand): $8 \times 2^3 = 64$.\item Two separate blocks of 2 overlapping pairs, or a chain of 3 pairs plus 1 isolated pair (forces 6 people to stand): $36 \times 2^2 = 144$.\item One block of 2 overlapping pairs plus 2 isolated pairs (forces 7 people to stand): $24 \times 2^1 = 48$.\item Four completely disjoint pairs (forces all 8 people to stand): There are exactly $2$ ways to space out four separate pairs around an 8-seat table ($e_1, e_3, e_5, e_7$ or $e_2, e_4, e_6, e_8$). This leaves 0 people free: $2 \times 2^0 = 2$.
                \end{itemize}


                Summing these gives $64 + 144 + 48 + 2 = 258$.

 \textbf{Step 5: Intersections of five pairs}There are $\dbinom{8}{5} = 56$ choices. Grouping by the number of unique people forced to stand:

 
                \begin{itemize}
                    \item Chain of 5 pairs (forces 6 people): $8 \times 2^2 = 32$.\item Mixed configurations forcing 7 people: $32 \times 2^1 = 64$.\item Mixed configurations forcing all 8 people: $16 \times 2^0 = 16$.
                \end{itemize}


                Summing these gives $32 + 64 + 16 = 112$.

 \textbf{Step 6: Intersections of six pairs}There are $\dbinom{8}{6} = 28$ choices:

 
                \begin{itemize}
                    \item Chain of 6 pairs (forces 7 people): $8 \times 2^1 = 16$.\item Configurations forcing all 8 people: $20 \times 2^0 = 20$.
                \end{itemize}


                Summing these gives $16 + 20 = 36$.

 \textbf{Step 7: Intersections of seven pairs}There are $\dbinom{8}{7} = 8$ choices. Any combination of 7 pairs forms an open chain that wraps almost completely around the table, meaning all $8$ people are forced to stand. 
 \[8 \times 2^0 = 8\]

 \textbf{Step 8: Intersection of all eight pairs}There is $\dbinom{8}{8} = 1$ choice, which forms the complete closed circle and forces all 8 people to stand.
 \[1 \times 2^0 = 1\]

 \textbf{Piecing it together:}Using the alternating signs required by PIE, the total count of distinct bad configurations is:
 \[\text{Bad Configurations} = 512 - 576 + 448 - 258 + 112 - 36 + 8 - 1 = 209\]

 Subtracting these invalid outcomes from our initial total pool yields the number of successful valid configurations:
 \[\text{Valid Configurations} = 256 - 209 = 47\]

 Thus, the probability that no two adjacent people stand up is $\boxed{\textbf{(A) } \dfrac{47}{256}}$.

 


	\item (\textit{AMC 10 10A 2015}) The zeroes of the function $f(x)=x^2-ax+2a$ are integers. What is the sum of the possible values of $a$?


 $\textbf{(A) }7\qquad\textbf{(B) }8\qquad\textbf{(C) }16\qquad\textbf{(D) }17\qquad\textbf{(E) }18$


 
 \textbf{Answer: }C



\textbf{Solution 1}

By Vieta's Formula, $a$ is the sum of the integral zeros of the function, and so $a$ is integral.

 Because the zeros are integral, the discriminant of the function, $a^2 - 8a$, is a perfect square, say $k^2$. Then adding 16 to both sides and completing the square yields
 \[(a - 4)^2 = k^2 + 16.\]Therefore $(a-4)^2 - k^2 = 16$ and
 \[((a-4) - k)((a-4) + k) = 16.\]Let $(a-4) - k = u$ and $(a-4) + k = v$; then, $a-4 = \dfrac{u+v}{2}$ and so $a = \dfrac{u+v}{2} + 4$. Listing all possible $(u, v)$ pairs (not counting transpositions because this does not affect ($u + v$)), $(2, 8), (4, 4), (-2, -8), (-4, -4)$, yields $a = 9, 8, -1, 0$. These $a$ sum to $16$, so our answer is $\boxed{\textbf{(C) }16}$.

 



\textbf{Solution 2}

Let $r_1$ and $r_2$ be the integer zeroes of the quadratic. By \href{https://artofproblemsolving.com/wiki/index.php?title=Vieta\%27s\_Formulas}{Vieta's Formulas}, 
 \[r_1 + r_2 = a\text{ and }r_1r_2 = 2a.\]

 Plugging the first equation in the second, 
 \[r_1r_2 = 2 (r_1 + r_2).\] 

 Rearranging gives 
 \[r_1r_2 - 2r_1 - 2r_2 = 0 \implies (r_1 - 2)(r_2 - 2) = 4.\]

 These factors $(f_1,f_2)$ (ignoring order, because we want the sum of factors), can be $(1, 4), (-1, -4), (2, 2),$ or $(-2, -2)$.

 The sum of distinct $a = r_1 + r_2 = (f_1+2) + (f_2+2)$, and these factors give 
 \[\sum_a a = (5+4) + (-5+4) + (4+4) + (-4+4) = \boxed{\textbf{(C) }16}\].

 


 



\textbf{Solution 3}

Let the roots be $r$ and $s$. By Vieta's we have that $r+s = a$, and $rs = 2a$. So by Simon's Favorite Factoring Trick: 
 \[2r+2s = rs \Longrightarrow rs-2r-2s = 0 \Longrightarrow (r-2)(s-2) = 4.\] Now, we test out values of $r$ and $s$ to see which work. Lets say that $r-2 = 1$, and $s-2 = 4$. This implies that $(r,s) = (3,6)$, so $a = 9$. Now, we let $r-2 = 2$, and $s-2 = 2$, this gives us $(r,s) = (4,4)$, which gives us a value for $a$ of $8$. Now, we circle an answer of $17$ and move onto the next problem. Yikes! They didn't say that $a$ can't be negative! This is where many people would go wrong. Let $r-2 = -1$, and $s-2 = -4$. This gives us that $(r,s) = (1,-2)$. This results in a value of $-1$ for $a$. Now, we let $r-2 = -2$, and $s-2 = -2$. This gives $a = 0$. Now, we are done. Our total sum is $9+8-1 = \boxed{\textbf{(C) }16}$.

 -jb2015007 - minor edits by mathlover1205 and PojoDotCom

 



\textbf{Solution 4: Quadratic Formula}

We are given that the zeroes are integers. For this reason, we should use the quadratic formula on $x^2 - ax + 2a$ to yield more information about the roots. 

 After simplifying, we get that the roots are $\dfrac{a \pm \sqrt{a^2 - 8a}}{2}$. From this, we can see that a\textasciicircum{}2 - 8a has to be a perfect square for the roots to be integers. 

 In addition, $a^2 - 8a$ and $a$ have to have the same parity (both even or both odd). However, that is always true (you can try out some odd/even integers)

 Hence, I set up another equation: $a^2 - 8a = y^2$, where y is an integer.

 I use the quadratic formula again for $a$, yielding that the roots are $4 \pm \sqrt{y^2 + 16}$. If we sum up 2 conjugate roots, the sum is $8$, so the problem boils down to finding nonnegative "Pythagorean Triples" with $16$. 

 Hence, we get that $y = 3$, $y = 0$. Each contributes $8$ to the running sum, so our answer is $8 + 8 = 16$.

 Note: We can see that there are at max 2 values for $y$, as the answer choices indicate either $8$ or $16$.

 
P.S. Thanks for the latex formatting. ~DARSHSHETTY4

 


	\item (\textit{AMC 10 10A 2015}) For some positive integers $p$, there is a quadrilateral $ABCD$ with positive integer side lengths, perimeter $p$, right angles at $B$ and $C$, $AB=2$, and $CD=AD$.  How many different values of $p<2015$ are possible?


 $\textbf{(A) }30\qquad\textbf{(B) }31\qquad\textbf{(C) }61\qquad\textbf{(D) }62\qquad\textbf{(E) }63$


 
 \textbf{Answer: }B



\textbf{Solution 1}

Let $BC = x$ and $CD = AD = y$ be positive integers. Drop a perpendicular from $A$ to $CD$ to show that, using the Pythagorean Theorem, that
 \[x^2 + (y - 2)^2 = y^2.\]Simplifying yields $x^2 - 4y + 4 = 0$, so $x^2 = 4(y - 1)$. Thus, $y$ is one more than a perfect square.

 The perimeter $p = 2 + x + 2y = 2y + 2\sqrt{y - 1} + 2$ must be less than 2015. Simple calculations demonstrate that $y = 31^2 + 1 = 962$ is valid, but $y = 32^2 + 1 = 1025$ is not. On the lower side, $y = 1$ does not work (because $x > 0$), but $y = 1^2 + 1$ does work. Hence, there are 31 valid $y$ (all $y$ such that $y = n^2 + 1$ for $1 \le n \le 31$), and so our answer is $\boxed{\textbf{(B) } 31}$

 



\textbf{Solution 2}

Let $BC = x$ and $CD = AD = z$ be positive integers. Drop a perpendicular from $A$ to $CD$. Denote the intersection point of the perpendicular and $CD$ as $E$.

 $AE$'s length is $x$, as well.Call $ED$ $y$.By the Pythagorean Theorem,  $x^2 + y^2 = (y + 2)^2$.And so: $x^2 = 4y + 4$, or $y = (x^2-4)/4$.

 Writing this down and testing, it appears that this holds for all $x$. However, since there is a dividend of 4, the numerator must be divisible by 4 to conform to the criteria that the side lengths are positive integers.In effect, $x$ must be a multiple of 2 to let the side lengths be integers. We test, and soon reach 62. It gives us $p = 1988$, which is less than 2015. However, 64 gives us $2116 > 2015$, so we know 62 is the largest we can go up to. Count all the even numbers from 2 to 62, and we get $\boxed{\textbf{(B) } 31}$.

 -jackshi2006

 



\textbf{Solution 3 (Less Counting)}

Let $BC = n$ and $CD = AD = m$ be positive integers. Drop a perpendicular from $A$ to $CD$. Denote the intersection point of the perpendicular and $CD$ as $E$.

 $AE$'s length is $n$, and $ED = m-2$By the Pythagorean Theorem,  $m^2 = (m-2)^2 + n^2$, therefore $n^2 = 4(m-1)$

 Notice that 4 is already a perfect square and $n$ is an integer, meaning $m-1$ also has to be a perfect square. We can denote $m-1$ as $x^2$, therefore we have $m = x^2 + 1$, $n = 2x$, and each $x$ will correspond to a pair of $(n, m)$. 

 Since $p<2015$, $2 + n + m + m<2015$ which gives us $n+2m<2013$. Substituting $x$ will give us $2x+2(x^2 + 1) < 2013$, simplify this will again give us $x^2 + x + 1 < 1006.5$Try $x = 32$, $x^2 + x + 1 = 1057 > 1006.5$, exceededso $x$ must be less than $32$

 When $x=31$, $x^2 + x + 1 = 993 < 1006.5$, $x=31$ works for us.

 For $x$ in $(0, \infty)$, $f(x) = x^2 + x + 1$ increases, therefore for any $x<31$, $f(x) < f(31) < 1006.5$

 Thus, any $x$ in $[1, 31]$ satisfy the condition that $x^2 + x + 1 < 1006.5$. $x$ needs to be integer, and there are $31$ integers in the interval $[1, 31]$. 

 Each $x$ corresponds to a $(n, m)$, meaning there are $31$ possible value for different pairs of $(n, m)$ and lead to $\boxed{\textbf{(B) } 31}$ possible values for the perimeter $p$. 

 


	\item (\textit{AMC 10 10A 2015}) Let $S$ be a square of side length $1$.  Two points are chosen at random on the sides of $S$.  The probability that the straight-line distance between the points is at least $\frac12$ is $\frac{a-b\pi}c$, where $a$, $b$, and $c$ are positive integers with $\gcd(a,b,c)=1$.  What is $a+b+c$?


 $\textbf{(A) }59\qquad\textbf{(B) }60\qquad\textbf{(C) }61\qquad\textbf{(D) }62\qquad\textbf{(E) }63$


 
 \textbf{Answer: }A



\textbf{Solution 1 (Calculus)}

Divide the boundary of the square into halves, thereby forming $8$ segments. Without loss of generality, let the first point $A$ be in the bottom-left segment. Then, it is easy to see that any point in the $5$ segments not bordering the bottom-left segment will be distance at least $\dfrac{1}{2}$ apart from $A$. Now, consider choosing the second point on the bottom-right segment. The probability for it to be distance at least $0.5$ apart from $A$ is $\dfrac{0 + 1}{2} = \dfrac{1}{2}$ because of linearity of the given probability. (Alternatively, one can set up a coordinate system and use geometric probability.)

 If the second point $B$ is on the left-bottom segment, then if $A$ is 

 


\end{enumerate}